\documentclass[8pt]{beamer}
\usetheme{Madrid}
\usefonttheme{serif}

\usepackage{fontspec}
% надёжный шрифт с кириллицей (работает в Overleaf)
\setmainfont{Libertinus Serif}[Ligatures=TeX]
\setsansfont{Libertinus Sans}[Ligatures=TeX]
\newfontfamily\cyrillicfont{Libertinus Serif}
\newfontfamily\cyrillicfontsf{Libertinus Sans}
\usepackage{polyglossia}
\setmainlanguage{russian}
\setotherlanguage{english}

\usepackage{amsmath}
\usepackage{amsfonts}
\usepackage{amssymb}
\usepackage{graphicx}
\usepackage{microtype}

\usepackage{tabularx}
\usepackage{booktabs}

\usepackage{xcolor}
\usepackage{pifont}
\usepackage{tikz}
\newcommand{\plus}{\tikz[baseline=(X.base)] \node[fill=green!65!black, rounded corners=2pt, inner sep=2pt] (X) {\color{white}\ding{51}};}
\newcommand{\minus}{\tikz[baseline=(X.base)] \node[fill=red!70!black,   rounded corners=2pt, inner sep=2pt] (X) {\color{white}\ding{55}};}

\DeclareMathOperator{\sen}{sen}
% \DeclareMathOperator{\tg}{tg}

\setbeamertemplate{caption}[numbered]

% --- Информация о презентации ---
\author[Автор]{Чиняев Владимир Николаевич}
\title{Практические аспекты орбитальных операций \\ Анализ ошибок}
\newcommand{\coursename}{Введение в маневрирование КА}
\setbeamertemplate{navigation symbols}{} 
\institute[]{Московский физико технический институт} 
\date{2 декабря 2025 г.} 

\bibliographystyle{apalike}

% --- Кастомная нижняя синяя полоска: курс слева, номер слайда справа ---
\setbeamertemplate{footline}{
  \leavevmode%
  \begin{beamercolorbox}[wd=\paperwidth,ht=2ex,dp=1ex]{palette primary}
    \hspace{0.4cm}\usebeamerfont{footline}\coursename\hfill\usebeamerfont{footline}\insertframenumber\hspace{0.4cm}
  \end{beamercolorbox}%
}

% Чтобы убрать лишние отступы в footer-шрифте (опционально)
\setbeamerfont{footline}{size=\small}

\begin{document}

\begin{frame}
\titlepage
\end{frame}

\begin{frame}{Анализ ошибок}
Всё наше предшествующее обсуждение маневрирования предполагало, что манёвры выполняются точно в соответствии с планом. К сожалению, в реальных миссиях \textbf{манёвры никогда не выполняются так, как они были запланированы}. \\
\smallskip
\textbf{Основные источники отклонений:}
\begin{enumerate}
    \item Даже самые лучшие модели прогноза баллистического движения - это лишь модели реального движения;
    \item Существуют неопределённости в каждом измерении, которое производится в ходе миссии. Небольшие ошибки в определении положения и ориентации приведут к отклонению от номинального плана;
    \item Довольно часто непредвиденные изменения влияют на ход миссии. Например, потеря контакта с наземной станцией, отказ ПО, ошибка передачи данных и тд;
    \item КА ведёт себя не так, как было рассчитано. Первый манёвр в последовательности обычно сильно отличается от рассчитанного;
    \item Невозможно полностью исключить человеческие ошибки.
\end{enumerate}

\textbf{Каждый из этих факторов способен привести к существенным отклонениям в выполнении миссии.}
\end{frame}

\begin{frame}{Анализ ошибок в жизненном цикле миссии}
Предположим, что мы планируем предстоящую миссию. 
\begin{itemize}
    \item Первый шаг - рассчитать номинальный план маневрирования КА. На этом этапе используются все те алгоритмы, что мы изучали ранее в ходе курса;
    \item Далее, рассчитанное маневрирование подвергается всеобъемлющему тестированию для того, чтобы убедиться в устойчивости решения к неопределённостям и ошибкам в планировании и исполнении манёвров.
\end{itemize}
Такой анализ проводится до и после запуска:
\begin{itemize}
    \item \textbf{До запуска} - чтобы убедиться в том, что проект маневрирования не имеет серьёзных проблем;
    \item \textbf{После запуска} - чтобы учесть реальные данные и состояние КА, убедиться в надёжности миссии.
\end{itemize}

Мы можем представить катастрофическую поведение КА в ходе миссии, требующее радикальных изменений в структуре миссии, включающих:
\begin{itemize}
    \item Изменение ПО;
    \item Изменение орбиты;
    \item Изменение методов поддержания орбиты;
    \item и тд.
\end{itemize}
Далее мы ограничим рассмотрение сценариями, в которых характеристики КА находятся в разумных пределах, относительно номинальных значений. 
\end{frame}

\begin{frame}{Главные последствия неноминальной производительности}
\begin{enumerate}
    \item \textbf{Потери топлива}
    \begin{itemize}
        \item Номинальная последовательность манёвров оптимизируется по топливу (или по $\Delta V$), что напрямую определяет максимальный срок службы КА на орбите;
        \item Любое отклонение от номинальной траектории приводит к неоптимальному расходу топлива;
        \item Поскольку манёвры никогда не выполняются идеально, отклонения от номинального плана неизбежны.
    \end{itemize}

    \item \textbf{Потери времени}
    \begin{itemize}
        \item Время выполнения маневрирования также критично, так как определяет загрузку персонала и наземной инфраструктуры;
        \item Пропуск манёвра приводит к добавлению как минимум одного витка к длительности миссии; при наличии фазовых ограничений может потребоваться несколько дополнительных витков.
    \end{itemize}
\end{enumerate}
\end{frame}

\begin{frame}{Классификация рассматриваемых ошибок}
Мы обсудим 5 типов анализа ошибок:
\begin{enumerate}
    \item \textbf{Анализ ошибок вывода КА ракетой-носителем} \\
    Показывает, что целевая орбита может быть достигнута, если ракета-носитель отклонилась от номинальной траектории в заданном интервале;

    \item \textbf{Анализ ошибок наведения} \\
    Показывает, что цели миссии могут быть достигнуты, если направление импульсов отличается от запланированного;

    \item \textbf{Анализ опозданий манёвров} \\
    Показывает, что цели миссии могут быть достигнуты, если манёвр был задержан на небольшое количество времени (менее витка);

    \item \textbf{Анализ пропущенных манёвров} \\
    Показывает, что цели миссии могут быть достигнуты, если манёвр был задержан более чем на виток;

    \item \textbf{Анализ ошибок ДУ} \\
    Показывает, что цели миссии могут быть достигнуты, если величина импульсов отличается от запланированного.
    
\end{enumerate}
\end{frame}

\begin{frame}{Ошибки вывода КА ракетой-носителем}
Идеально, если ракета-носитель каждый раз доставляет КА точно на орбиту выведения. В реальности, конечно, ракета-носитель может достичь орбиты выведения в заданных пределах точности, продиктованных неопределённостями в 
\begin{itemize}
    \item Моделировании,
    \item Работе ДУ,
    \item Работе датчиков,
    \item и тд.
\end{itemize}
Дисперсия может быть задокументирована как набор отклонений орбитальных параметров или в виде матрицы ковариации. Вне зависимости от специфической формы, разработчик КА должен быть уверен, что целевая орбита будет достигнута вне зависимости от орбиты выведения.

Для анализа дисперсии орбиты выведения берётся набор возмущённых орбит выведения (в диапазоне $3 \sigma$) и выполняется полный расчёт и оптимизация маневрирования. После анализа всех возмущённых случаев максимальное количество дополнительного топлива выделяется как резервное.

Существуют два основных способа выбора потенциальных орбит выведения:
\begin{enumerate}
    \item \textbf{Метод Монте-Карло} \\
    Случайно сгенерировать большое число орбит выведения из предоставленных данных о распределении и для каждой рассчитать оптимальное маневрирование;

    \item \textbf{Граничные случаи} \\
    Просто протестировать граничные случаи для значений орбитальных элементов, предполагая, что именно они и потребуют максимальное количество дополнительного топлива. 
\end{enumerate}
\end{frame}

\begin{frame}{Ошибки наведения}
Для того чтобы приложить импульс $\Delta V$ один (или более) двигателей КА должны быть сонаправлены с направлением $\Delta V$. Это требует изменения ориентации КА. Такие изменения обычно производятся с помощью маховиков, гиродинов или двигателей ориентации. Вне зависимости от метода, целевая ориентация не будет достигнута с идеальной точностью. Стоит различать:
\begin{enumerate}
    \item \textbf{Целевая ориентация} \\
    Ориентация, в которой двигатели точно совпадают с направлением импульса;

    \item \textbf{Истинная ориентация} \\
    Ориентация, которая реально достигается КА в подготовке к приложению импульса;

    \item \textbf{Измеренная ориентация} \\
    Ориентация, которую имеет оператор, основываясь на данных измерений датчиков (солннечный датчик, звёздный датчик, акселерометр и тд.)
\end{enumerate}

Поскольку истинная ориентация отличается от целевой, то импульс $\Delta V$ будет приложен в неоптимальном направлении. Если миссия протекает нормально, то отклонение между этими тремя величинами будет небольшим. 

Маховики и гиродины являются более точными в отличие от двигателей ориентации. Можно привести следующие грубые оценки:
\begin{itemize}
    \item \textbf{Маховики и гиродины:} \\
    Могут достичь целевой ориентации с точностью $0.1^\circ$ ($1^\circ$ в случае отсутствия калибровки);

    \item \textbf{Двигатели ориентации:} \\
    Точность сильно зависит от минимальной продолжительности импульса, точности датчиков, величины переориентации, направления переориентации и тд. Точность будет составлять от $1^\circ$ (номинально) до $5^\circ$ ($3 \sigma$).
\end{itemize}
\end{frame}

\begin{frame}{Ошибки наведения}
Поскольку ещё до выполнения переориентации известно, что фактическая ориентация неизбежно будет отличаться от требуемой, заранее выполняют анализ чувствительности манёвров к ошибкам наведения. Проверяется каждый манёвр, чтобы убедиться, что вся последовательность остаётся устойчивой к этим отклонениям. \\
\smallskip
Для каждого манёвра оценивается расход топлива в наихудшем случае ошибки ориентации. Поскольку реализация всех худших случаев одновременно маловероятна, обычно используют корень из суммы квадратов этих затрат и включают полученную величину в запас топлива как резерв на ошибки наведения. \\
\smallskip
В процессе моделирования и оптимизации при анализе ошибок наведения используется следующий набор правил:
\begin{itemize}
    \item Время начала манёвра — номинальное;
    \item Длительность манёвра — номинальная;
    \item Характеристики работы ДУ (тяга и удельный импульс) — номинальные;
    \item Орбита КА известна точно.
\end{itemize}
\end{frame}

\begin{frame}{Ошибки наведения}
\textbf{После манёвра с ошибкой наведения следует ли вставлять в последовательность перед следующим манёвром небольшой корректирующий манёвр переориентации?} \\
\smallskip
\begin{itemize}
    \item Дополнительная переориентация действительно увеличивает степень свободы и может дать небольшую экономию топлива.
    \begin{itemize}
        \item Эта экономия часто несущественна, если ошибка наведения мала или у аппарата есть достаточный топливный резерв;
        \item Любой дополнительный манёвр несёт риск, поэтому если допустимая траектория существует без переориентации, её обычно выбирают как базовую.
    \end{itemize}

    \item Команда управления миссией должна за ограниченное время сформировать номинальный сценарий и провести необходимые проверки. Часто приоритеты смещены в сторону других задач, не связанных с анализом ошибок. Поэтому в рамках анализа ошибок наведения допустимо использовать последовательность манёвров, которая является «достаточно хорошей», а не строго оптимальной.
\end{itemize}
\end{frame}

\begin{frame}{Ошибки наведения}
\textbf{Должна ли оставшаяся часть последовательности манёвров строго следовать номинальному плану или же баллистикам разрешается по своему усмотрению добавлять, переносить или удалять манёвры?} \\
\smallskip
\begin{itemize}
    \item \textbf{Гибкий подход} \\ 
    Для переходной фазы заранее выделяется временное окно, в рамках которого аналитикам орбиты разрешено свободно добавлять, изменять или удалять манёвры, если это помогает достичь целевой орбиты. Такой подход встречается в программах с надёжным и удобным доступом к собственной сети наземных станций.

    \item \textbf{Жёсткий подход} \\
    Последовательность должна максимально следовать согласованному заранее номинальному плану. Новые манёвры не добавляются (и существующие не исключаются), если только это не критически необходимо.  \\
    Основной аргумент — фактическая траектория остаётся в окрестности заранее проанализированных сценариев, для которых был выполнен детальный анализ ошибок до запуска; отклонения считаются «хорошо изученными».
\end{itemize}
\end{frame}

\begin{frame}{Предполётная оценка ошибок эффективности ДУ}
Двигатели КА проходят огневые испытания на Земле. На испытательном стенде измеряется создаваемая тяга и расход компонентов топлива. На основе этих данных инженер по ДУ определяет \textbf{удельный импульс ДУ}.

Характеристики работы двигателя и расхода топлива и окислителя в полёте, как правило, не совпадает точно с результатами наземных испытаний. Различия между характеристиками на Земле и в полёте могут быть объяснены
\begin{itemize}
    \item изменением давления в баках под воздействием тепловой обстановки в космосе;
    \item небольшими отличиями в расходных характеристиках магистралей подачи компонентов;
    \item погрешностями моделей;
    \item тем фактом, что невозможно испытать двигатель на Земле во всех условиях, соответствующих эксплуатационным режимам на орбите.
\end{itemize}
До запуска необходимо оценить наихудший расход топлива на случай, если характеристики ДУ на орбите окажутся значительно ниже, чем измеренные на Земле.
\end{frame}

\begin{frame}{Предполётная оценка ошибок тяги ДУ}
ДУ имеет отклонения и по величине создаваемой тяге. Фактический уровень тяги может оказаться ниже номинального. Планировщики миссии не будут знать о том, что ДУ работает на пониженном уровне тяги до выполнения первого манёвра, а следовательно, спланируют всю последовательность исходя из номинального уровня. Уменьшение тяги приведёт к тому, что последующие манёвры необходимо планировать с большей продолжительностью, чтобы
\begin{itemize}
    \item реализовать исходно запланированное $\Delta V$;
    \item компенсировать недобор $\Delta V$ при первом манёвре.
\end{itemize}
Удлинение манёвров может оказаться проблемой, если существуют ограничения на длительность манёвров, связанные, например, с тепловыми или энергетическими условиями. Ещё одно последствие более длительных манёвров - увеличение потерь эффективности при переходе от мгновенных манёвров к протяжённым.
\end{frame}

\begin{frame}{Предполётная оценка последствий отказа ДУ}
Распространённой схемой дизайна КА, которому необходимо сильно изменить характеристики орбиты, является установка одного высокоэффективного маршевого двигателя. Предполагается, что именно эта ДУ выполнит весь переход, после чего КА перейдёт в режим штатной эксплуатации. В дальнейшем такая ДУ не используется или используется крайне редко. Часто основной двигатель изолируется от остальной ДУ для избегания его непреднамеренного запуска. \\
\smallskip
В целом такие маршевые двигатели обладают высокой надёжностью; однако они не застрахованы от отказов. \\
\textbf{Возможны различные проблемы:}
\begin{itemize}
    \item Неправильно смонтированная электрическая цепь питания приводит к короткому замыканию, из-за чего двигатель невозможно повторно зажечь;
    \item Ошибка в бортовом программном обеспечении вызывает немедленное прерывание манёвра;
    \item Неправильно установленный клапан подачи рабочего тела блокирует его поступление в двигатель.
\end{itemize}
В одних ситуациях отказ проявляется уже при первом включении, полностью выводя двигатель из эксплуатации; в других — двигатель может успешно работать в течение лет, пока не возникнет неисправность.
\end{frame}

\begin{frame}{Предполётная оценка последствий отказа ДУ}
При необходимости перехода к другой комбинации двигателей или к другой двигательной установке в целом могут возникать дополнительные затраты рабочего тела за счёт четырёх различных эффектов:
\begin{enumerate}
    \item Меньший удельный импульс непосредственно приводит к большему расходу топлива;
    \item Более длительные манёвры приводят к большим потерям от преобразования импульсных манёвров в протяжённые;
    \item Если последовательность манёвров растягивается, удовлетворить ограничения миссии может стать сложнее, что потенциально увеличивает расход рабочего тела;
    \item Если последовательность перелёта манёвров, орбитальные возмущения могут оказывать заметное неблагоприятное влияние на орбиту, что также может привести к увеличению расхода рабочего тела.
\end{enumerate}

Важно отметить, что резервный комплект двигателей может использовать полностью отдельный запас рабочего тела. В этом случае, если ни один из доступных двигателей не может использовать рабочее тело, связанное с отказавшим двигателем, этот запас превращается в бесполезную массу. Один из примеров — когда основная ДУ использует пару топливо/окислитель, а резервная ДУ использует только топливо.
\end{frame}

\begin{frame}{Повторимость импульсов ДУ}
Предыдущие вопросы были посвящены анализу различных сценариев до запуска. Такие сценарии рассматриваются по двум причинам:
\begin{enumerate}
    \item Наличие заранее подготовленной информации позволяет быстро принимать взвешенные решения при возникновении нештатной ситуации;
    \item Позволяет оценить необходимый запас топлива и при необходимости учесть его.
\end{enumerate}

По сути, такие ошибки можно рассматривать как внесение систематического смещения в работу ДУ. \\
\smallskip
Например, в случае пониженной эффективности предполагается, что удельный импульс снижен (смещён) до значения $3 \sigma$ ниже номинала на протяжении всей последовательности манёвров. \\
\smallskip
Существует и другой эффект, типично проявляющийся в ходе реальных миссий: разброс характеристик двигателя от манёвра к манёвру, называемый \textbf{повторяемостью работы двигателя} (thruster repeatability). Он проявляется в том, что фактические параметры близки к прогнозируемым, но не совпадают с ними точно.
\end{frame}

\begin{frame}{Повторимость импульсов ДУ}
Повторяемость характеристик работы двигателя является сложной функцией множества параметров, среди которых два наиболее важных — давление и температура ДУ. Специалисты обычно используют детальные модели двигательной системы для двух задач:
\begin{enumerate}
    \item прогнозирования температур, давлений и эффективности перед манёвром;
    \item восстановления фактической работы двигателя по телеметрии после манёвра.
\end{enumerate}
Такие модели подробно описывают все элементы: 
\begin{itemize}
    \item баки,
    \item трубопроводы,
    \item клапаны,
    \item регуляторы,
    \item двигатели и их конфигурацию — вплоть до всех изгибов труб.
\end{itemize}
Однако даже при тщательном моделировании предсказать работу двигателя бывает трудно, особенно для первого экземпляра новой серии аппаратов.
\end{frame}

\begin{frame}{Повторимость импульсов ДУ}
Для конкретного КА существует множество факторов, которые могут приводить к изменению характеристик работы двигателя от одного манёвра к другому. Краткий (но не исчерпывающий) перечень включает следующие аспекты:
\begin{itemize}
    \item Разная ориентация аппарата во время манёвров меняет освещённость и температурный профиль, что вызывает небольшие изменения в работе двигательной установки;
    \item Если между манёврами проходит значительное время (например, сутки), условия входа и выхода из тени, длительность пребывания в тени могут поменяться, что также влияет на температурный режим КА;
    \item Температурная обстановка зависит от работы обогревателей — с прошлого манёвра какие-то нагреватели могли быть включены или отключены;
    \item В системах с регулированием давления входное давление двигателя зависит от точки запирания регулятора, которая может незначительно меняться; даже малые колебания влияют на параметры тяги;
    \item За время между манёврами немного изменяется относительное положение Солнца, что также изменяет тепловую среду аппарата.
\end{itemize}
\textbf{Все эти факторы не могут быть учтены точно}, поэтому существуют некоторые вариации в реальной работе ДУ в отличие от прогнозируемой.
\end{frame}

\begin{frame}{Неопределённость в соотношении компонентов топливной смеси}
\textbf{Ошибки в соотношении компонентов топливной смеси} приводят к тому, что один из компонентов оказывается в избытке. Этот избыток не сгорает и превращается в бесполезную массу, фактически увеличивая топливные затраты на манёвры. \\
\smallskip
\textbf{Можно выделить два источника проблемы:}
\begin{itemize}
    \item Ошибка в заправочном количестве компонентов. Возникает при неверной оценке требуемого смесевого отношения или прогнозируемой производительности ДУ.
    \item Отклонение фактического расхода компонентов в двигателе от номинального. Возможные причины: частичное засорение или пережатие трубопровода, сбой работы регулятора расхода и т. п.
\end{itemize}
\end{frame}

\begin{frame}{Опоздание манёвра}
\begin{center}
    \textbf{Способы выполнения манёвров}
\end{center}

\begin{columns}
    \begin{column}{0.50\textwidth}
        \begin{center}
            \textbf{По таймеру}
        \end{center}
        \plus\ КА может выполнять манёвры вне зависимости от видимости с наземной станции; \\
        \minus\ Операторы не могут вмешаться в работу КА в случае сбоя.
    \end{column}

    \begin{column}{0.50\textwidth}
        \begin{center}
            \textbf{По команде ЦУП}
        \end{center}
        \plus\ Операторы могут непосредственно отслеживать состояние КА в момент исполнения манёвра; \\
        \minus\ Существует риск пропустить точное время старта из-за сбоев связи (помехи, погода, проблемы оборудования).
    \end{column}
\end{columns}
\smallskip
Если момент выполнения манёвра упущен, необходимо сформировать новую последовательность манёвров. Однако нередко проблему удаётся устранить быстро, и возможность передачи команды восстанавливается за считанные минуты. В такой ситуации может возникнуть вопрос, возможно ли выполнить манёвр позже запланированного времени? Как сильно мы можем задержать начало манёвра?
\end{frame}

\begin{frame}{Опоздание манёвра}
Чтобы подготовиться к такой ситуации можно выполнить \textbf{анализ позднего включения двигателя}. Фиксируются параметры манёвра (ориентация, длительность, тяга), а время начала последовательно сдвигается позже. Для каждого варианта оставшаяся последовательность переоптимизируется, и строится кривая топливного штрафа за задержку включения. \\
\smallskip
Если манёвр пропущен, можно просто отказаться от него и сформировать новую номинальную последовательность, добавив целые витки до следующего окна для манёвра. Однако, это увеличивает время перехода и повышает операционные затраты. Поэтому иногда выгоднее выполнить манёвр с небольшой задержкой, даже если это вызывает небольшой перерасход топлива. \\
\smallskip
\textbf{Разные миссии по-разному подходят к учёту таких ситуаций:}
\begin{enumerate}
    \item \textbf{Задать максимальную допустимую задержку выполнения манёвра} (например, $10$ минут). После этого анализируются все манёвры в последовательности и вычисляется максимальный возможный топливный штраф для каждого из них. В топливный бюджет добавляется отдельная статья, равная этому худшему случаю. Здесь предполагается, что \textbf{возможна только одна задержка манёвра}.

    \item \textbf{Выделить фиксированное количество рабочего тела на задержку манёвра в топливном бюджете.} Далее каждый манёвр анализируется индивидуально, чтобы определить, насколько поздно его можно выполнить, прежде чем будет израсходован этот выделенный топливный резерв. В таком подходе допустимая задержка будет различаться для разных манёвров (например, первому допустимо $5$ минут, второму — $10$ минут и т.д.).
\end{enumerate}
\end{frame}

\begin{frame}{Пропуск манёвра}
Ранее мы оценивали потери топлива в ситуации, когда манёвр всё же выполняется, но с задержкой на несколько минут. Начиная с некоторого момента задержка становится настолько затратной по топливу, что владелец/оператор предпочитает полностью пропустить данное окно для манёвра и выполнить его в следующую доступную возможность — на последующем витке. Это и называется \textbf{пропущенный манёвр}. \\
\smallskip
\textbf{Основные проблемы:}
\begin{enumerate}
    \item \textbf{Иной набор орбитальных возмущений} \\
    Аппарат должен был быть на новой орбите; вместо этого он остаётся на старой и испытывает другой набор возмущений.

    \item \textbf{Активные ограничения} \\
    Пропуск одного манёвра сдвигает все последующие манёвры на виток-два вперёд.
    \begin{itemize}
        \item относительное положение Солнца меняется, следовательно, учёт некоторых ограничений может вызывать проблемы.
        \item расход топлива может заметно увеличиться.
    \end{itemize}

    \item \textbf{Нарушение фазирования} \\
    Если требуется выйти на заданную фазу в группировке, пропуск манёвра нарушает фазу.
\end{enumerate}
\end{frame}

\begin{frame}{Аварийное прекращение манёвра}
Ещё одна возможная ситуация — манёвр может быть начат, но затем прекращён до запланированного времени окончания. Это называется \textbf{аварийным прекращением манёвра}. Прерывание может произойти в любой момент: через доли секунды после начала, за секунду до окончания или в любой промежуточный момент. \\
\smallskip
\textbf{Причин аварийного прекращения может быть много:}
\begin{itemize}
    \item Превышение температурных ограничений двигателя;
    \item Выход скорости вращения аппарата за допустимые пределы;
    \item Нарушение ограничений по углу на Солнце;
    \item Ошибка ПО;
    \item Ошибка оператора и тд.
\end{itemize}
После прерывания манёвра космический аппарат оказывается на \textbf{незапланированной} орбите:
\begin{itemize}
    \item это уже не исходная (до манёвра) орбита;
    \item но и не целевая (послеманевровая) орбита.
\end{itemize}
Прежде чем разбирать план восстановления после прерванного манёвра, необходимо решить более неотложную задачу. Наземным станциям слежения требуется \textbf{первоначальная оценка положения КА}, чтобы начать поиск и захват сигнала.
\end{frame}

\begin{frame}{Аварийное прекращение манёвра}
Возможно также выполнить ограниченный \textbf{анализ ошибок для сценария прерванного манёвра}. Но
\begin{itemize}
    \item Заранее неизвестно, в какой именно момент манёвр может быть прерван;
    \item Выполнять полноценную оптимизацию миссии для каждой возможной секунды манёвра — слишком трудоёмко по времени и вычислительным ресурсам.
\end{itemize}
Простым решением является выбор нескольких характерных точек вдоль длительности манёвра. \\
\smallskip
\textbf{Такой анализ даёт три ключевых преимущества:}
\begin{itemize}
    \item Подтверждается выполнимость миссии — по крайней мере, для нескольких репрезентативных точек вдоль длительности манёвра.
    \item Оцениваются возможные потери топлива, вызванные прерыванием манёвра.
    \item Уточняется понимание возможных вариантов действий, а также трудностей планирования в различных сценариях аварийного завершения импульса.
\end{itemize}
Очевидно, что \textbf{безопасность космического аппарата имеет абсолютный приоритет после прерванного манёвра}. Все вопросы восстановления орбиты после аварийного прерывания рассматриваются только после того, как аппарат будет возвращён в безопасное состояние.
\end{frame}

\begin{frame}[allowframebreaks]{Источники}
  \begin{thebibliography}{9}

    \bibitem{DiPrinzio}
    DiPrinzio M. Methods of Orbital Maneuvering. – American Institute of Aeronautics and Astronautics, Inc., 2019.
    
  \end{thebibliography}
\end{frame}

\end{document}
